% =====================================================================
\section{Projekt i implementacja rozwiązania}
\label{sec:impl}

\subsection{Stopień realizacji wymagań}

Prototyp PlayLink osiągnął poziom \textbf{wersji zaawansowanej ($\geq$90\%)}. Wszystkie
15 wymagań funkcjonalnych z tab.~\ref{sec:wf} oraz wszystkie 12 wymagań niefunkcjonalnych
z tab.~\ref{sec:wnf} zostały zrealizowane i wdrożone na środowisku produkcyjnym
(\url{https://playlink.bartek.monster}). Aplikacja jest w pełni działająca:
uwierzytelnianie BIP39, tworzenie/dołączanie/opuszczanie pokojów, czat realtime z
podpisami, kalendarz i RSVP, wyszukiwanie i filtrowanie, panel administratora oraz
automatyczne czyszczenie wygasłych pokojów.

Zastosowano aktualne standardy projektowe: typowanie statyczne w całym stosie
(TypeScript + Pydantic/SQLModel), automatyczną walidację żądań, migracje bazy
(Alembic) oraz w pełni zautomatyzowany pipeline CI/CD (rozdz.~\ref{sec:zarz}).

\subsection{Projekt graficznego interfejsu użytkownika}
\label{sec:gui}

Interfejs zaprojektowano w estetyce inspirowanej przeglądarką lobby z gry
\textit{Diablo~II: Resurrected} — ciemna, gotycka paleta złoto-brązowa i
minimalistyczny układ. Frontend zbudowano w SvelteKit~2 / Svelte~5 (runes) z Tailwind
CSS; UI reaguje natychmiast na zdarzenia WebSocket. Poniżej wybrane widoki aplikacji.

\begin{figure}[H]
  \centering
  \begin{subfigure}[t]{0.48\textwidth}
    \includegraphics[width=\textwidth]{prez-slide25.png}
    \caption{Strona startowa („Enter realm”).}
  \end{subfigure}\hfill
  \begin{subfigure}[t]{0.48\textwidth}
    \includegraphics[width=\textwidth]{d2-rooms-v1.png}
    \caption{Przeglądarka pokojów (lista lobby).}
  \end{subfigure}

  \vspace{0.7em}
  \begin{subfigure}[t]{0.48\textwidth}
    \includegraphics[width=\textwidth]{smoke-form-filled.png}
    \caption{Formularz tworzenia pokoju.}
  \end{subfigure}\hfill
  \begin{subfigure}[t]{0.48\textwidth}
    \includegraphics[width=\textwidth]{smoke-detail-with-metadata.png}
    \caption{Szczegóły pokoju z metadanymi.}
  \end{subfigure}

  \vspace{0.7em}
  \begin{subfigure}[t]{0.48\textwidth}
    \includegraphics[width=\textwidth]{prez-slide29.png}
    \caption{Widok pokoju: uczestnicy, wydarzenie z RSVP, czat.}
  \end{subfigure}\hfill
  \begin{subfigure}[t]{0.48\textwidth}
    \includegraphics[width=\textwidth]{prez-slide35.png}
    \caption{Moderacja: zdarzenie systemowe o wyrzuceniu uczestnika.}
  \end{subfigure}
  \caption{Wybrane widoki graficznego interfejsu PlayLink.}
  \label{fig:gui}
\end{figure}

\subsection{Aspekty AI/ML}
\label{sec:aiml}

W uzgodnieniu z opiekunem projektu (mgr inż. Bartłomiej Jencz) zespół świadomie
zrezygnował z komponentu AI/ML na rzecz \textbf{rozszerzonego, dopracowanego modelu
bezpieczeństwa danych} (rozdz.~\ref{sec:bezp}). Decyzja wynika z charakteru produktu:
kluczową wartością PlayLink jest natychmiastowy, bezpieczny i prywatny dostęp
(tożsamość niepowiernicza BIP39, weryfikowalne autorstwo wiadomości), a nie
rekomendacje oparte na uczeniu maszynowym. Dobór pingu odbywa się deterministycznie
na podstawie odległości geograficznej lobby (metryka Haversine), bez modelu ML.

\subsection{Aspekty bezpieczeństwa danych}
\label{sec:bezp}

Bezpieczeństwo jest centralnym filarem PlayLink. Przyjęto zasadę \emph{non-custodial}:
klucz prywatny nigdy nie opuszcza przeglądarki, a backend wyłącznie weryfikuje dowody
kryptograficzne. Poniżej uzgodnione aspekty bezpieczeństwa — wszystkie zrealizowane
(pokrycie $\geq$90\%).

\paragraph{1. Tożsamość bezhasłowa (BIP39 + ECDSA).}
Konto reprezentuje 12-słowna fraza mnemoniczna BIP39 (standard z Bitcoina, 2013;
2048-wyrazowy słownik daje 132 bity entropii). Z frazy derywowany jest klucz
secp256k1, a publiczny adres EIP-55 stanowi identyfikator konta. Eliminuje to e-mail,
numer telefonu i hasło — a tym samym wektor ataku „reset hasła”.

\paragraph{2. Uwierzytelnianie challenge-response i ochrona przed replay.}
Logowanie wykorzystuje jednorazowy \emph{nonce}: backend wydaje wyzwanie, klient
podpisuje je (EIP-191 \code{personal_sign}), a backend odtwarza adres z podpisu
(ECDSA recovery). Zabezpieczenia:
\begin{itemize}
  \item nonce jest \textbf{jednorazowy} (flaga \code{used}) i \textbf{krótkożyciowy} (TTL 5~min);
  \item w bazie przechowywany jest wyłącznie \textbf{skrót SHA-256} wartości nonce —
        surowa wartość nigdy nie trafia na dysk;
  \item ponowne wydanie nonce unieważnia poprzedni — co blokuje ataki powtórzeniowe.
\end{itemize}

\paragraph{3. Izolacja tokenu sesji (wzorzec BFF, hardening XSS).}
Token JWT (HS256) jest przechowywany w ciasteczku \code{httpOnly},
\code{sameSite=strict}, ustawianym przez serwer SvelteKit. Token \textbf{nigdy nie jest
dostępny dla JavaScriptu klienta}, co znacząco ogranicza skutki ewentualnego XSS.

\paragraph{4. Weryfikowalne autorstwo wiadomości.}
Wiadomości czatu są podpisywane kluczem derywowanym z BIP39 (EIP-191). Serwer
rekonstruuje kanoniczną postać wiadomości i weryfikuje podpis, a znaczniki czasu
dryfujące ponad \code{CHAT_SIGNATURE_SKEW_SECONDS} są odrzucane — co ogranicza
powtórzenia. Autorstwo jest weryfikowalne \emph{end-to-end}, a nie tylko na podstawie JWT.

\paragraph{5. Autoryzacja administratora bez eskalacji.}
Uprawnienia administratora wynikają wyłącznie z listy adresów w zmiennej środowiskowej
\code{ADMIN_ADDRESSES} — w bazie nie istnieje kolumna \code{is_admin}. Uniemożliwia to
eskalację uprawnień przez modyfikację danych; sfałszowane roszczenie o admina jest
odrzucane (zob.\ testy bezpieczeństwa, rozdz.~\ref{sec:testy}).

\paragraph{6. Ograniczanie nadużyć (rate limiting) i higiena treści.}
REST jest limitowany przez slowapi (domyślnie 10/min), a czat — przesuwnym oknem
per-adres. Edycja nazwy użytkownika przechodzi przez filtr wulgaryzmów (stoplist).

\paragraph{7. Bezpieczeństwo sieci i wdrożenia.}
Cały ruch szyfrowany jest TLS~1.3 (HTTPS + WSS, certyfikaty Let's Encrypt). Wdrożenie
jest self-hosted: transfer artefaktów odbywa się prywatnym kanałem przez Tailscale
(mesh VPN, WireGuard), publiczny SSH nie jest wystawiony (brak ataków brute-force na
port~22), a baza PostgreSQL nie ma mapowania portu na hosta.

\paragraph{Weryfikacja.}
Wszystkie powyższe mechanizmy są pokryte automatycznymi testami bezpieczeństwa
(replay nonce, wygasły nonce, sfałszowane roszczenie admina, zmanipulowane podpisy) —
zob.\ rozdz.~\ref{sec:testy}.
